\chapter{面向多类别有害知识的分级可控编辑框架}\label{chap:experiment}

\section{相关知识介绍}
本节为示例占位内容。表标题和表序号居中置于表上方，如表~\ref{tab:params} 所示。

\begin{table}[htbp]
  \centering
  \caption{示例表格——主要参数设置}
  \label{tab:params}
  \begin{tabular}{lcc}
    \toprule
    参数 & 符号 & 取值 \\
    \midrule
    参数一 & $\alpha$ & $0.1$ \\
    参数二 & $\beta$  & $0.9$ \\
    参数三 & $N$      & $1000$ \\
    \bottomrule
  \end{tabular}
\end{table}

\subsection{过度拒绝问题}
% TODO: 一刀切编辑的现象与例证，安全-效用失衡分析。

\subsection{分级编辑形式化定义}
% TODO: 编辑强度作为内容风险函数的形式化定义；分级映射可替换设计。

\subsection{因果追踪技术}
% TODO: activation patching / causal tracing 原理与流程（为 3.2 探测做预备）。

\section{方案框架}

\subsection{探测方法设计}
% TODO: 探测方法设计：暴力、违法建议、自残、偏见、越狱五类危害类别的逐层探测。

\subsection{层级分布图谱}
% TODO: 产出「危害类别 × 模型层」分布图谱；高低危分布差异检验——分级依据的实证发现。

\subsection{连续强度调度}
% TODO: α(x) = sigmoid(f_risk(x))；L_edit = α·L_forget + (1-α)·L_retain。

\subsection{差异化编辑策略}
% TODO: 低危方向性衰减（保留知识本体）；高危零空间擦除（借鉴 AlphaEdit 零干扰方向）。

\subsection{分级锚定机制}
% TODO: 衔接第 2 章锚定机制：高危强锚定、低危弱锚定，CKEdit 框架在此汇合。

\section{实验结果与分析}
本节为示例占位内容。不同方法的对比结果如表~\ref{tab:results} 所示，其中表头
使用了跨行单元格。

\begin{table}[htbp]
  \centering
  \caption{示例表格——结果对比}
  \label{tab:results}
  \begin{tabular}{lccc}
    \toprule
    \multirow{2}{*}{方法} & 指标一 & 指标二 & 指标三 \\
                          & (\%)   & (\%)   & (s)    \\
    \midrule
    方法 A     & $86.0$ & $5$ & $42.3$ \\
    方法 B     & $91.5$ & $3$ & $38.7$ \\
    本文方法   & $\mathbf{98.2}$ & $\mathbf{0}$ & $\mathbf{31.5}$ \\
    \bottomrule
  \end{tabular}
\end{table}

由定理~\ref{thm:demo} 与式~\eqref{eq:demo1} 可知，此处可进一步展开对结果的分析说明。

\subsection{实验设置}
% TODO: 数据集与分级映射（SafeEdit 九类映射三级）；模型；基线（ToxEdit、统一强度）。

\subsection{安全效用平衡分析}
% TODO: E1/E2/E4：攻击成功率 + 过度拒绝率 + 分层安全效用 Pareto 曲线。

\subsection{分级精度与泄漏分析}
% TODO: 两项自建指标——编辑是否恰好只影响目标风险等级；是否波及其他等级。

\subsection{基线方法对比}
% TODO: E3 ToxEdit 三线对比：toxicity 正面战 / 多类别扩展战 / 推理开销对比。

\subsection{消融与敏感性分析}
% TODO: E5 统一中强度；E6 离散三档 vs 连续调度；E7 分级映射替换敏感性。

\section{本章小结}
本节为示例占位内容。请替换为本章的小结。
